\section{Architecture comparison with OpenRocket}\label{app:or-architecture}

The two preceding appendices compared what the tools offer (\cref{app:or-features}) and how
their physics is computed (\cref{app:or-methods}). This appendix compares them as software:
how each is built and modularized, how each represents a rocket in memory, how each
positions parts and computes mass properties, how each structures a simulation run, where
each lets outside code in, and how each persists its work. The framing of
\cref{app:or-context} applies throughout --- QtRocket is an independent project, OpenRocket
is a calibration baseline, and the comparison measures distance, not conformance. It is also
worth saying once what the comparison is \emph{for}: the two architectures were shaped by
different pressures (a seventeen-year-old product with a large user base; a young core built
seams-first), and reading them side by side shows which of QtRocket's choices are genuine
alternatives and which are simply earlier points on the same road.

Provenance differs by side, as everywhere in these appendices. OpenRocket statements were
researched from the project's repository at tag \code{release-24.12} and branch
\code{unstable} (examined 2026-07-07) and from its published developer
documentation\footnote{\url{https://github.com/openrocket/openrocket} and
\url{https://openrocket.readthedocs.io/en/latest/dev_guide/architecture.html}. Class, enum,
and package names quoted below are exact; short code excerpts are verbatim from the tagged
sources.}; they are not file-and-line verified the way QtRocket claims are. QtRocket
statements carry \code{file:line} citations into the pinned commit \code{254cd41}.
\Cref{tab:or-architecture-summary} is the map; the subsections walk it row by row.

\begin{table}[htbp]
\centering
\small
\begin{tabularx}{\textwidth}{@{}lXX@{}}
\toprule
\textbf{Dimension} & \textbf{OpenRocket 24.12} & \textbf{QtRocket \code{254cd41}} \\
\midrule
Language / build & Java 17, Gradle, two JPMS modules & C++23, CMake, layered static libraries \\
\addlinespace
Wiring & Guice injector behind a static locator & Controller singleton; constructor injection at the composition root \\
\addlinespace
Design tree & 23 concrete \code{RocketComponent} classes, per-class child-compatibility rules & 5 \code{Part} types in one composite, 4 factory-constructible \\
\addlinespace
Mass properties & External \code{MassCalculator}, recomputed per step, hookable & In-tree composite cache behind dirty flags and a mass-delta gate \\
\addlinespace
Positioning & \code{AxialMethod} + \code{RadiusMethod} offsets; nose-tip origin, $+x$ aft & \code{StationLink} placement intent + one resolver; nose-tip datum, \zfwd \\
\addlinespace
Simulation control & Event queue, 16 event types, four phase steppers, forking per stage & Single loop, typed termination, no events \\
\addlinespace
Extensibility & Per-step listener hooks, GraalJS scripting, annotation-scanned plugins & No plugin surface; machine-scriptable REPL protocol \\
\addlinespace
Persistence & Versioned \code{.ork} ZIP+XML; byte-signature multi-format import & Single-format \code{.qrd} XML; placement intent; fail-closed loads \\
\addlinespace
Motor data & Bundled snapshot (SQLite pipeline on \code{unstable}) & Ships empty; three ingest paths including live REST \\
\addlinespace
GUI & Mature Swing product with live background simulation & Thin Qt front end; authoring is CLI-first \\
\bottomrule
\end{tabularx}
\caption{The architecture comparison at a glance. Each row is expanded, with sources, in the
subsection that follows it.}\label{tab:or-architecture-summary}
\end{table}

\subsection{Language, runtime, and build}\label{app:or-architecture-build}

OpenRocket is Java: \code{Application.java} hard-codes
\cppinline|SUPPORTED_JRE_VERSIONS = { 17 }|, and CI builds against Temurin~17. The build is
Gradle, producing two subprojects that are also two Java Platform Module System (JPMS)
modules, each with its own \code{module-info.java}: \code{info.openrocket.core} --- the
simulation engine, file formats, and design classes, deliberately headless and reusable ---
and \code{info.openrocket.swing}, the GUI, which depends on core and carries the application
entry point. Object wiring uses Google Guice~7: a static service locator,
\code{info.openrocket.core.startup.Application}, holds the injector, and code obtains
services through it (\code{FlightConfiguration}, for example, calls
\cppinline|Application.getInjector().getInstance(RocketDescriptor.class)|); the older
generation of static accessors is marked deprecated in favor of the injector
(``\code{@deprecated} Fetch the db from Guice instead'').

QtRocket is C++23 \src{CMakeLists.txt}{5} under CMake, and its modularity is the link graph
rather than a module system: \code{utils} links only third-party math and networking
\src{utils/CMakeLists.txt}{14}, \code{sim} links only \code{utils}
\src{sim/CMakeLists.txt}{28}, \code{model} links \code{utils} and \code{sim}
\src{model/CMakeLists.txt}{44}, and the controller compiles once into the
\code{qtrocket\_core} static library that every executable links \src{CMakeLists.txt}{224}.
Qt is confined to the GUI executable and the visualizer; the core, like OpenRocket's, is
headless by construction (\cref{sec:overview}). Wiring is manual: the controller is a
process-wide singleton \src{QtRocket.cpp}{19} acting as the composition root, and the one
dependency that must vary per run --- the \code{Environment} --- is injected into the
\code{Propagator}'s constructor there \src{QtRocket.cpp}{48}, never reached through a
global.

The honest version of the ``DI container versus singletons'' contrast is narrower than the
labels suggest. OpenRocket's injector is itself resolved through a static locator, so most
call sites still reach a global --- what the container buys is \emph{rebindability}: tests
and plugins can install different bindings without touching call sites, and the plugin
framework (\cref{app:or-architecture-ext}) depends on exactly that. QtRocket's object graph
is small enough to wire by hand; it pays for its two singletons with a known blemish (the
controller's double-checked initialization reads its flag outside the mutex
\src{QtRocket.cpp}{21}, where the \code{Logger} one layer down already uses the safe Meyers
form \src{utils/Logger.cpp}{19}) and gets its rebindability from ordinary constructor
seams instead --- the injected \code{Environment}, and the fake-injectable
\code{ThrustCurveAPI} interface behind the motor database \src{model/ThrustCurveClient.h}{68}.

\begin{designrule}[One instinct, two enforcement mechanisms]
Both projects enforce the same first-order rule --- a headless core that must not know
about its GUI. OpenRocket enforces it with the module system: \code{swing} requires
\code{core}, never the reverse, and the JVM checks the boundary. QtRocket enforces it with
the link graph and one comment stated where it binds: ``Bottom layer: utils must not depend
on model/ or sim/'' \src{utils/CMakeLists.txt}{13}, with Qt appearing only in front-end link
lists \src{CMakeLists.txt}{286}. A module system checks the rule mechanically; a link graph
makes violating it a build failure. Either way, the invariant is load-bearing, not
documentation.
\end{designrule}

\subsection{The design tree and mass properties}\label{app:or-architecture-tree}

Both projects model a rocket as a tree of components; the trees differ in breadth and,
more interestingly, in where the rules live. OpenRocket's
\code{info.openrocket.core.rocketcomponent} package defines roughly 23 concrete classes
under three abstract families --- \code{ComponentAssembly} (the \code{Rocket} root, stages,
pods), \code{ExternalComponent} (body tubes, transitions, nose cones, four fin-set classes,
lugs, rail buttons), and \code{InternalComponent} (rings, couplers, bulkheads, mass objects,
recovery devices) --- with cross-cutting capability interfaces such as \code{MotorMount} and
\code{RingInstanceable}. Tree legality is enforced per class: every component implements
\cppinline|isCompatible(Class<? extends RocketComponent> type)|, and \code{addChild()}
throws when the answer is no. \code{Rocket} accepts only \code{AxialStage}; \code{BodyTube}
accepts pods, internals, and any external component that is not itself a body --- fins
attach to tubes, tubes do not nest in tubes. The type system carries the domain rules.

QtRocket's composite has five concrete \code{Part} types registered in one umbrella header
\src{model/parts/Parts.h}{4}, four of them constructible through a factory that dispatches
on the type name and validates required fields \src{model/parts/PartFactory.cpp}{27} over a
single type-agnostic parameter struct \src{model/parts/PartFactory.h}{27} (the \code{Motor}
is deliberately excluded \src{model/parts/PartFactory.cpp}{73}). Attachment rules are
structural rather than typed: \code{addChildPart} rejects null, already-parented, and
cycle-forming children \src{model/parts/Part.cpp}{71}, but any part may parent any other ---
there is no \code{isCompatible} analogue. What QtRocket has instead is a \emph{geometric}
legality gate: the placement sweep of \cref{sec:placement} checks the resolved tree for
self-intersection and fails the mass computation closed when it finds one
\src{model/parts/Part.cpp}{184}. OpenRocket asks ``may this class contain that class?'' at
edit time; QtRocket asks ``does this geometry physically fit?'' at resolve time. The two
checks are complementary, and neither project currently has both.

The mass-properties split is the genuinely different philosophy. OpenRocket keeps mass
\emph{out} of the components: the \code{masscalc} package is deliberately small ---
\code{MassCalculator}, \code{MassCalculation}, and \code{RigidBody}, the aggregate
mass-and-inertia result --- and the calculator walks the component tree afresh whenever the
simulation needs mass properties; nothing is cached inside components. QtRocket computes
composite mass, CM, and inertia \emph{inside} the tree: a two-pass, time-aware walk
\src{model/parts/Part.cpp}{175} whose per-part tensors follow the per-unit-mass convention
\src{model/parts/Part.h}{30} and whose results are cached behind a structural dirty flag and
the mass-delta gate \src{model/parts/Part.cpp}{137} (\cref{sec:parts}).

\begin{designrule}[OpenRocket: components describe, a calculator computes]
Because mass computation is a free-standing service rather than a component method, it is
naturally an interception point: the listener interface of
\cref{app:or-architecture-ext} exposes \code{preMassCalculation}/\code{postMassCalculation}
hooks that can replace or adjust the \code{RigidBody} a step sees. The components stay a
pure description of the rocket; every consumer --- simulator, analysis dialogs, plugins ---
goes through one calculator. The price is recomputation on every step and no natural home
for a cache.
\end{designrule}

\begin{keyinsight}[QtRocket: the cache that can never be wrong --- under one assumption]
Storing the composite in the tree buys QtRocket a property OpenRocket's per-step calculator
does not have: after burnout, when \code{getMass(t)} repeats bit-for-bit, the cache freezes
and repeated or rejected integrator stage queries are served bit-stably
\src{model/parts/Part.h}{248}, while the ungated live mass sum stays cheap as the ODE
divisor \src{model/parts/Part.cpp}{118}. The price is that the gate keys on \emph{total}
mass alone, so its soundness is coupled to the single-motor assumption
\src{model/RocketModel.cpp}{13} --- two simultaneous burners holding total mass constant
while shifting CM would be served stale properties. Neither design dominates: OpenRocket
pays per step for generality; QtRocket pays in coupled assumptions for bit-stability.
\end{keyinsight}

\subsection{Positioning models}\label{app:or-architecture-pos}

OpenRocket positions components with a pair of enums. Each component stores an
\code{(AxialMethod, offset)} pair, where \code{AxialMethod} is one of \code{ABSOLUTE},
\code{AFTER}, \code{TOP}, \code{MIDDLE}, \code{BOTTOM}; each value implements
\cppinline|getAsPosition(offset, innerLength, outerLength)| returning the component's
front-face coordinate in the parent's frame (\code{MIDDLE}, for instance, is
$\mathtt{offset} + (\mathtt{outerLength} - \mathtt{innerLength})/2$). \code{AFTER} is the
internal sibling-stacking mode for body components in sequence; the GUI exposes the other
four. Radial placement is a second enum, \code{RadiusMethod} (\code{COAXIAL}, \code{FREE},
\code{RELATIVE}, \code{SURFACE} --- the last a flush mount on the parent's outer surface).
The arithmetic is purely geometric: component lengths only, never mass or CM. The global
frame puts the origin at the nose tip with $+x$ pointing \emph{aft}.

QtRocket stores a \code{StationLink} per attachment --- fractional stations on both ends of
the seam, a gap, and a seat kind \src{model/parts/Placement.h}{49}, with the seat kind
closed at three values because a rigid axisymmetric stack admits only three radial
relationships \src{model/parts/Placement.h}{32} --- and derives every absolute pose through
one resolver, ``intent (StationLink) in, geometry (Pose) out''
\src{model/parts/Placement.h}{148}, whose core is two lines of arithmetic
\src{model/parts/Placement.cpp}{20}. The frame puts the datum at the nose tip with \zfwd:
each part occupies $z \in [-L, 0]$ in its own frame \src{model/parts/Part.h}{130}. The
frame-convention difference is therefore concrete: both projects plant the origin at the
nose tip, but their longitudinal axes run in \emph{opposite directions} --- a station aft
of the tip has positive $x$ in OpenRocket and negative $z$ in QtRocket, and OpenRocket uses
$x$ where QtRocket uses $z$. Any future \code{.ork} import will live or die on getting that
sign flip right once, in one place.

The deeper distinction is what is stored. OpenRocket stores geometry directly: the offset
is a signed length between anchor planes, computed by the author, and the method names
which planes. QtRocket stores intent and derives the geometry: the link names which feature
of the child seats against which feature of the parent, the seat kind selects the gap's
sign and the radial compatibility rule, and the resolver --- pure geometry, ``No CM, time, or
mass'' \src{model/parts/Placement.h}{158} --- turns it into coordinates on
every read (\cref{sec:placement}). The gap between the two is smaller than it first appears:
OpenRocket's \code{TOP}/\code{MIDDLE}/\code{BOTTOM} anchors are themselves re-evaluated
against live lengths at query time, so both systems survive a parent-length edit. What
OpenRocket's scheme does not carry is radial semantics in the axial store (radius is a
separate method with no compatibility meaning) or a validation gate coupled to placement;
what QtRocket's scheme does not carry is OpenRocket's \code{FREE}/\code{RELATIVE} radial
freedom --- off-axis seating is deliberately deferred to 6-DOF (\cref{sec:incomplete}).
Both projects converged independently on the same load-bearing rule: placement arithmetic
must never involve mass. QtRocket arrived there the hard way, by replacing a CM-coupled
legacy scheme (\cref{sec:placement}); OpenRocket appears to have started there.

\subsection{The simulation engine}\label{app:or-architecture-engine}

This is the largest structural difference between the two codebases. OpenRocket's engine,
\code{BasicEventSimulationEngine}, is event-driven: a time-ordered \code{EventQueue} rides
on the mutable \code{SimulationStatus}, the run begins by pushing a \code{LAUNCH} event, and
the loop alternates between handling due events and advancing the ODE, clamping each step so
it never overshoots the next queued event. The \code{FlightEvent.Type} enum names sixteen
event kinds --- among them \code{IGNITION}, \code{LIFTOFF}, \code{LAUNCHROD},
\code{BURNOUT}, \code{EJECTION\_CHARGE}, \code{STAGE\_SEPARATION}, \code{APOGEE},
\code{RECOVERY\_DEVICE\_DEPLOYMENT}, \code{GROUND\_HIT}, \code{TUMBLE}, and
\code{SIMULATION\_END}. Integration is delegated to a \code{SimulationStepper} chosen per
flight phase (\cref{lst:or-architecture-steppers}): full 6-DOF RK4 under power and coast, a
simplified 3-DOF stepper under a deployed recovery device, a tumble stepper, and a ground
stepper after landing. Stage separation \emph{forks the simulation}: the engine pushes a
copy of \code{SimulationStatus} onto a \code{toSimulate} stack and later pops it to fly the
separated booster, each branch producing its own \code{FlightDataBranch} in the result.

\begin{listing}[htbp]
\begin{minted}{java}
private final SimulationStepper flightStepper  = new RK4SimulationStepper();
private final SimulationStepper landingStepper = new BasicLandingStepper();
private final SimulationStepper tumbleStepper  = new BasicTumbleStepper();
private final SimulationStepper groundStepper  = new GroundStepper();
\end{minted}
\caption{The four phase steppers of OpenRocket's event-driven engine, quoted from
\code{BasicEventSimulationEngine.java} at tag \code{release-24.12}. The engine swaps
steppers as events change the flight phase; the fidelity of the dynamics is a per-phase
decision, not a global one.}\label{lst:or-architecture-steppers}
\end{listing}

QtRocket's engine is a single loop: \code{Propagator::runUntilTerminate()}
\src{sim/Propagator.cpp}{58} advances one solver (RK4 or the embedded-error RKF45, selected
at runtime through the integrator facade \src{sim/Integrator.h}{52}) until the model's own
predicate says the flight is over --- descending below the launch site
\src{model/RocketModel.cpp}{62} --- or one of the guard exits fires. There is no event
queue, no phase steppers, and no forking; what QtRocket has instead is the typed
five-value termination taxonomy \src{sim/Propagator.h}{35} that partitions every possible
exit into a diagnosis (\cref{sec:simcore}). The two designs are honest reflections of their
feature sets: OpenRocket needs an event architecture because recovery, staging, and
tumbling \emph{are} events; QtRocket's every flight is single-phase and ballistic, so a
loop with a predicate is the whole requirement. The event system is precisely the
architectural piece QtRocket's planned recovery and staging work will need, and the
termination taxonomy is its natural seed (\cref{sec:incomplete}).

\subsection{Extension surfaces}\label{app:or-architecture-ext}

OpenRocket lets outside code into the running simulation. Three listener interfaces attach
to a run: \code{SimulationListener} (start/end/pre-step/post-step, with a boolean pre-step
veto), \code{SimulationEventListener} (observe or veto flight events), and --- the key
mechanism --- \code{SimulationComputationListener}, which brackets \emph{every physics
sub-computation} with paired hooks: \code{pre}/\code{post} around acceleration, the
atmospheric model, wind, gravity, flight conditions, aerodynamics, mass calculation, and
thrust. A \code{pre*} hook returning non-null replaces the model's output for that step; a
\code{post*} hook may modify the result. On top of the listeners sit simulation extensions
(the tree ships a \code{JavaCode} extension and script-backed ones), JavaScript scripting through GraalJS behind the
standard JSR-223 API, and a plugin framework in which providers are discovered by classpath
annotation scan (\code{@Plugin}, ClassGraph) and bound through Guice modules --- runtime
plugin JARs, enabled by the container architecture of
\cref{app:or-architecture-build}.

QtRocket has no listener, plugin, or scripting mechanism of any kind. Its pluggability is
internal and compile-time-shaped: string-keyed registries select among the built-in
gravity, atmosphere \src{sim/Environment.h}{31}, and integrator implementations
\src{sim/Integrator.h}{52}, and the front ends enumerate those registries so their
vocabulary cannot drift from the core's \src{gui/SimOptionsTab.cpp}{32}. What QtRocket
offers outside code instead is the REPL: 33 commands over a machine-parseable stdout
protocol --- exactly one \code{OK}/\code{ERR} line per command, \code{CSV} paths, \code{\#}
comments \src{cli/Repl.cpp}{245} --- deliberately kept clean by silencing the logger to
ERROR in the CLI binary \src{cli/CliMain.cpp}{17} (\cref{sec:interfaces}). These are
different answers to different questions. OpenRocket extends \emph{inward}: injected code
runs between physics sub-steps and can replace a model mid-flight, which is what active
control experiments, custom wind fields, or air-starts need. QtRocket extends
\emph{outward}: an external program drives whole runs over a text protocol, which is what
parameter sweeps, CI harnesses, and reproducible batch studies need --- and which
OpenRocket itself has no interactive equivalent for (its headless story is a separate
Python helper module, \code{orhelper}\footnote{\url{https://github.com/openrocket/orhelper}
--- an org-maintained Python module ``to facilitate interacting and scripting with
OpenRocket.''}).
Neither surface subsumes the other: an outward protocol cannot touch a flight mid-step, and
an inward hook is no substitute for a scriptable process boundary.

\subsection{Persistence}\label{app:or-architecture-persist}

OpenRocket's design file, \code{.ork}, is a ZIP archive containing a versioned XML document
(\code{rocket.ork}) plus attached resources, and its format history is explicit: version
1.8 landed with release 22.02, 1.10 with 24.12, and 1.11 (embedding thrust curves in the
archive) is staged for the next release\footnote{Format-version table and archive layout:
\url{https://openrocket.readthedocs.io/en/latest/dev_guide/file_specification.html}.}.
The loader, \code{GeneralRocketLoader}, auto-detects what it is fed by byte signature ---
GZIP, ZIP, bare \code{<openrocket} XML, RockSim's \code{<RockSimDoc}, RASAero's
\code{<RASAeroDoc} --- so files from two foreign ecosystems and every historical
OpenRocket version load through one entry point; RockSim is supported for both import and
export. QtRocket persists designs in exactly one format: \code{.qrd} XML at version 0.2
\src{model/DesignSerializer.cpp}{220}, which stores placement \emph{intent} --- the
\code{<link>} elements --- and re-resolves geometry on every load
\src{model/DesignSerializer.cpp}{93}. Its loader is fail-safe in structure and fail-closed
in vocabulary: the tree is built detached and installed atomically so a bad file leaves the
current design untouched \src{model/DesignSerializer.cpp}{255}, an unknown seat kind is
rejected outright \src{model/DesignSerializer.cpp}{184}, and the retired 0.1 format is
refused with an explicit ``must be re-created'' rather than migrated
\src{model/DesignSerializer.cpp}{194} (\cref{sec:persistence}).

The postures could not be more different, and each is right for its context. OpenRocket's
loader is maximally forgiving because two decades of user files and three neighboring
ecosystems exist; rejecting any of them would strand real work. QtRocket's loader is
maximally strict because almost no files exist yet, so rejection is cheaper than a
migration framework and guarantees that anything loaded is coherent. The asymmetry has a
shelf life: the \code{.ork} version table is a precedent for exactly the machinery
\code{.qrd} will need once users accumulate files, and OpenRocket's content-sniffing entry
point is the natural shape for any future foreign-format import
(\cref{sec:incomplete}).

\subsection{Motor data}\label{app:or-architecture-motors}

OpenRocket ships its motor corpus with the application: release 24.12 bundles a 3.3\,MB
serialized snapshot of thrustcurve.org data (\code{thrustcurves.ser}), and the
\code{unstable} branch replaces it with an SQLite database
(\code{initial\_motors.db}, 1,458 motors and 1,591 curves per its metadata) published
from a dedicated data-pipeline repository\footnote{\url{https://github.com/openrocket/motor-database}
--- ``a motor database for OpenRocket based on thrustcurve.org data and other sources,''
republished as \code{motors.db.gz} for in-app updates.} --- motor data curation has been
moved out of the application into a data engineering pipeline. QtRocket inverts the
default: the shipped database is empty (the bundled \code{.rse} file is wired into test
binaries only, via a compile definition \src{tests/CMakeLists.txt}{16}), and motors arrive
at run time through three ingest paths behind one facade --- the RockSim \code{.rse} loader
\src{model/RSEDatabaseLoader.cpp}{24}, the strict validating RASP \code{.eng} parser
\src{model/RASPLoader.cpp}{137}, and a live thrustcurve.org REST client hidden behind the
fake-injectable \code{ThrustCurveAPI} seam \src{model/ThrustCurveClient.h}{68}. The
loaders are owned inside the database and never touched by front ends
\src{model/MotorModelDatabase.cpp}{63}; storage is a common-name-keyed map
\src{model/MotorModelDatabase.h}{129} (\cref{sec:motors}).

Architecturally this is offline-first versus online-first. OpenRocket's user is never
blocked by the network and never sees an empty chooser; the cost is a release-time curation
pipeline and data that ages between releases. QtRocket's user gets live search against the
canonical source --- a workflow the baseline's in-app selection does not lead with
(OpenRocket's core carries thrustcurve.org client classes, but its chooser is fed by the
bundled corpus) --- at the cost
of a first run that begins with nothing until a file is imported or the network answers.
The designs are converging from opposite ends: OpenRocket's new SQLite pipeline adds
in-app updating to a bundled corpus, and QtRocket's natural next step is bundling a
snapshot in front of its live client.

\subsection{GUI architecture}\label{app:or-architecture-gui}

OpenRocket's Swing GUI is the product: one document window hosting the 2-D rocket figure
and JOGL-based 3-D view, per-component configuration dialogs, plotting, printing, and an
optimizer. Its most architecturally interesting habit is \emph{live background simulation}:
\code{RocketPanel} owns a static \code{ExecutorService} --- ``a fixed-sized thread pool for
all background simulations with all threads in daemon mode and with minimum priority'' ---
and every design edit re-runs a simulation off the interaction thread to refresh the
altitude, velocity, and stability readouts drawn on the design canvas, alongside
continuously recomputed CP/CG markers. Simulation is treated as a rendering input, cheap
enough to run speculatively on every edit.

QtRocket's GUI is a thin front end over the same core its CLI drives, and at the pinned
commit it is the narrowest of the three surfaces: the design widgets are constructor-only
stubs \src{gui/RocketTreeView.cpp}{3}, the File actions ship disabled
\src{gui/MainWindow.ui}{117}, and design authoring lives in the REPL
(\code{newdesign} \src{cli/Repl.cpp}{784}, \code{addpart} \src{cli/Repl.cpp}{825}).
Its threading posture is the inverse of OpenRocket's twice over: the entire Qt event loop
runs on a \emph{spawned} \cppinline|std::thread| that constructs \code{QApplication}
\src{gui/GuiRunner.cpp}{62} --- an inversion of Qt's main-thread convention whose
consequences \cref{sec:interfaces} examines --- while simulation runs \emph{synchronously
on} that GUI thread from the launch button's slot \src{gui/CannonballTab.cpp}{116}.
OpenRocket pushes simulation off the interaction thread by default; QtRocket has not yet
separated the two at all. The one place the QtRocket rendering story is architecturally
ahead is outside the GUI proper: the standalone visualizer positions geometry with the
exact resolver the simulator's mass pass consumes (\cref{sec:interfaces}), a
single-source-of-truth property OpenRocket's separate figure-rendering path does not have.

\subsection{What each architecture makes easy}\label{app:or-architecture-close}

Read as wholes, the two architectures optimize for different verbs. OpenRocket's makes
\emph{extending the flight} easy: a new flight phenomenon is a new event type absorbed by
an engine already built around a queue; a new physics variant is a listener hook away; a
foreign file is one more byte signature; and the container wiring lets plugins rebind
services without touching call sites. The costs are the ones mature products accept ---
behavior assembled at run time by injection and annotation scan is harder to trace
statically, the loader's forgiveness dilutes invariants, and the simulation's
warn-and-continue posture (\cref{app:or-context}) means a produced number does not always
certify a coherent model. QtRocket's architecture makes \emph{trusting the number} easy:
the link graph is the layering, placement has one resolver shared by physics and rendering,
incoherent designs fail closed with located diagnostics, and a deterministic core behind a
text protocol makes every result reproducible from a script. Its costs are the mirror
image: no way into a flight mid-step, a closed persistence world, an empty first run, and
--- the expensive one --- no event architecture, so the recovery and staging features that
define the baseline's flight realism require growing a new control structure, not just new
physics (\cref{sec:incomplete}). Neither project would be improved by adopting the other's
architecture wholesale; each is shaped correctly for its stage. The useful reading of
OpenRocket for QtRocket's purposes is a preview of which seams get consumed next --- the
event queue, the background-simulation executor, the format-version table --- because the
baseline demonstrates, in shipped code, what those seams are for.
